\documentclass[a4paper,12pt]{article}

\usepackage[T2A]{fontenc}
\usepackage[utf8]{inputenc}
\usepackage[russian]{babel}
\usepackage{amsmath,amssymb,mathtools}
\usepackage{PTSans}
\renewcommand{\familydefault}{\sfdefault}
\usepackage{icomma}
\usepackage[
  left=20mm,
  right=20mm,
  top=19mm,
  bottom=21mm
]{geometry}
\usepackage{microtype}
\usepackage{xcolor}
\usepackage{tikz}
\usetikzlibrary{calc,arrows.meta,positioning,shapes.geometric}
\usepackage{tabularx,booktabs}
\usepackage{enumitem}
\usepackage{parskip}
\usepackage{fancyhdr}
\usepackage{setspace}
\usepackage{titlesec}

\onehalfspacing
\setlength{\parindent}{0pt}
\setlength{\parskip}{5pt}

% -------------------------------------------------
% Палитра
% -------------------------------------------------
\definecolor{accent}{HTML}{008080}
\definecolor{darkaccent}{HTML}{004D4D}
\definecolor{textgray}{HTML}{303638}
\definecolor{softgray}{HTML}{E8EEEE}

% -------------------------------------------------
% Служебные данные
% -------------------------------------------------
\newcommand{\student}{София Хоменко}
\newcommand{\teacher}{Лёвин Артём Александрович}

% -------------------------------------------------
% Колонтитулы
% -------------------------------------------------
\pagestyle{fancy}
\fancyhf{}
\fancyfoot[C]{\small\color{textgray}\thepage}
\renewcommand{\headrulewidth}{0pt}

% -------------------------------------------------
% Заголовки
% -------------------------------------------------
\titleformat{\section}
  {\Large\bfseries\color{darkaccent}}
  {}
  {0pt}
  {}

\titleformat{\subsection}
  {\large\bfseries\color{accent}}
  {}
  {0pt}
  {}

\titlespacing*{\section}{0pt}{1.1em}{0.45em}
\titlespacing*{\subsection}{0pt}{0.9em}{0.3em}

% -------------------------------------------------
% Типографические учебные блоки
% -------------------------------------------------
\newcommand{\keyidea}[1]{%
  \par\medskip
  {\bfseries\color{darkaccent}#1}\par
  \vspace{1pt}
  {\color{accent}\rule{\linewidth}{0.55pt}}\par
  \vspace{2pt}
}

\newcommand{\important}[1]{%
  \par\smallskip
  \noindent\textbf{\color{darkaccent}Важно.} #1\par
}

\newcommand{\exampletitle}{%
  \par\smallskip
  \noindent\textbf{\color{accent}Пример.}\enspace
}

% -------------------------------------------------
% Стили блок-схем
% -------------------------------------------------
\tikzset{
  flowstart/.style={
    ellipse,
    draw=darkaccent,
    line width=0.8pt,
    minimum width=43mm,
    minimum height=10mm,
    align=center,
    fill=softgray,
    font=\small
  },
  flowaction/.style={
    rectangle,
    rounded corners=2mm,
    draw=accent,
    line width=0.7pt,
    minimum width=70mm,
    minimum height=11mm,
    text width=66mm,
    align=center,
    fill=white,
    font=\small
  },
  flowside/.style={
    rectangle,
    rounded corners=2mm,
    draw=accent,
    line width=0.7pt,
    minimum width=50mm,
    minimum height=11mm,
    text width=46mm,
    align=center,
    fill=white,
    font=\small
  },
  flowdecision/.style={
    diamond,
    aspect=2.5,
    draw=darkaccent,
    line width=0.7pt,
    text width=39mm,
    align=center,
    inner sep=1.3mm,
    fill=softgray,
    font=\small
  },
  flowarrow/.style={
    -{Latex[length=3mm,width=2mm]},
    line width=0.8pt,
    draw=textgray
  },
  flowlabel/.style={
    font=\small,
    fill=white,
    inner sep=1pt
  }
}

\begin{document}

% =================================================
% ТИТУЛЬНЫЙ ЛИСТ
% =================================================
\thispagestyle{empty}

\begin{center}
\vspace*{20mm}

{\color{darkaccent}\Large\bfseries Чек-лист}

\vspace{6mm}

{\Huge\bfseries
Алгебраическое повторение\\[2mm]
и метод интервалов}

\vspace{8mm}

{\Large ОГЭ по математике}

\vspace{13mm}

{\large
Квадратные уравнения и разложение на множители\\
Алгебраические дроби и ограничения\\
Линейные неравенства и их системы\\
Системы линейных уравнений\\
Рациональные уравнения\\
Расчёты по формулам\\
Квадратные неравенства и метод интервалов}

\vfill

{\large\bfseries \teacher}

\vspace{2mm}

{\large эксклюзивно для \student}

\vspace{4mm}

{\large \today}

\vfill

\begin{tikzpicture}[x=1cm,y=1cm]
  \draw[accent,line width=1.1pt]
  (0,0)
  .. controls (3.0,0.75) and (8.3,-0.75) ..
  (13.3,0)
  .. controls (15.0,0.3) and (16.4,0.2) ..
  (16.8,0.05);
\end{tikzpicture}

\vspace{5mm}
\end{center}

\clearpage

% =================================================
% 1. КВАДРАТНЫЕ УРАВНЕНИЯ
% =================================================
\section{1. Квадратные уравнения и разложение на множители}

\keyidea{Основные формулы}
Квадратное уравнение имеет вид
\[
ax^2+bx+c=0,\qquad a\ne0.
\]

Дискриминант:
\[
D=b^2-4ac.
\]

Если \(D>0\), то
\[
x_{1,2}=\frac{-b\pm\sqrt D}{2a}.
\]
Если \(D=0\), то
\[
x=-\frac{b}{2a}.
\]
Если \(D<0\), действительных корней нет.

\important{Отрицательные коэффициенты подставляйте в формулу вместе со знаком. Например, если \(b=-1\), то в дискриминанте записываем \((-1)^2\).}

\exampletitle
Решим
\[
x^2-x-6=0.
\]
Здесь \(a=1\), \(b=-1\), \(c=-6\). Тогда
\[
D=(-1)^2-4\cdot1\cdot(-6)=25,
\]
\[
x_1=\frac{1+5}{2}=3,
\qquad
x_2=\frac{1-5}{2}=-2.
\]
Следовательно,
\[
x^2-x-6=(x-3)(x+2).
\]

\keyidea{Связь корней и множителей}
Если квадратный трёхчлен \(ax^2+bx+c\) имеет корни \(x_1\) и \(x_2\), то
\[
ax^2+bx+c=a(x-x_1)(x-x_2).
\]

\subsection{Неполные квадратные уравнения}
Не каждое квадратное уравнение удобно решать через дискриминант.

\exampletitle
\[
x^2-x=0
\quad\Longrightarrow\quad
x(x-1)=0
\quad\Longrightarrow\quad
x=0\ \text{или}\ x=1.
\]

\exampletitle
\[
x^2-6=0
\quad\Longrightarrow\quad
x^2=6
\quad\Longrightarrow\quad
x=\pm\sqrt6.
\]

\important{Из равенства \(x^2=a\) при \(a>0\) получаются два корня: \(x=\pm\sqrt a\).}

% =================================================
% 2. АЛГЕБРАИЧЕСКИЕ ДРОБИ
% =================================================
\section{2. Алгебраические дроби: ограничения и сокращение}

\keyidea{Главное правило}
Сокращать можно только \textbf{множители}. Поэтому перед сокращением числитель и знаменатель при необходимости раскладывают на множители.

\exampletitle
Упростим
\[
\frac{x^2-x-6}{x^2-4}.
\]
Сначала фиксируем ограничения исходной дроби:
\[
x^2-4\ne0
\quad\Longrightarrow\quad
x\ne-2,\qquad x\ne2.
\]
Разложим числитель и знаменатель:
\[
x^2-x-6=(x-3)(x+2),
\qquad
x^2-4=(x-2)(x+2).
\]
Тогда
\[
\frac{(x-3)(x+2)}{(x-2)(x+2)}
=
\frac{x-3}{x-2},
\qquad x\ne\pm2.
\]

\important{После сокращения ограничение \(x\ne-2\) сохраняется: исходное выражение при \(x=-2\) не определено.}

\keyidea{Формулы, которые стоит распознавать сразу}
\[
a^2-b^2=(a-b)(a+b),
\]
\[
(a+b)^2=a^2+2ab+b^2,
\qquad
(a-b)^2=a^2-2ab+b^2,
\]
\[
ax+ay=a(x+y).
\]

% =================================================
% FLOWCHART 1
% =================================================
\clearpage
\thispagestyle{plain}

\begin{center}
{\Large\bfseries\color{darkaccent}Алгоритм: упрощение алгебраической дроби}
\end{center}

\vspace{5mm}

\begin{center}
\begin{tikzpicture}[node distance=9mm]

\node[flowstart] (start) {Алгебраическая дробь};

\node[flowaction, below=of start] (restr)
{Записать ограничения: исходный знаменатель не равен нулю};

\node[flowaction, below=of restr] (factor)
{Разложить числитель и знаменатель на множители};

\node[flowdecision, below=12mm of factor] (common)
{Есть одинаковые множители?};

\node[flowside, below=17mm of common, xshift=-52mm] (cancel)
{Сократить одинаковые множители};

\node[flowside, below=17mm of common, xshift=52mm] (keep)
{Оставить дробь без сокращения};

\node[flowaction, below=24mm of common] (restore)
{Сохранить все ограничения исходного выражения};

\node[flowstart, below=of restore] (finish)
{Записать результат};

\draw[flowarrow] (start)--(restr);
\draw[flowarrow] (restr)--(factor);
\draw[flowarrow] (factor)--(common);

\draw[flowarrow]
  (common.west) -|
  node[flowlabel,pos=0.22]{да}
  (cancel.north);

\draw[flowarrow]
  (common.east) -|
  node[flowlabel,pos=0.22]{нет}
  (keep.north);

\draw[flowarrow] (cancel.south) |- (restore.west);
\draw[flowarrow] (keep.south) |- (restore.east);
\draw[flowarrow] (restore)--(finish);

\end{tikzpicture}
\end{center}

\clearpage

% =================================================
% 3. ЛИНЕЙНЫЕ НЕРАВЕНСТВА
% =================================================
\section{3. Линейные неравенства}

\keyidea{Алгоритм решения}
\begin{enumerate}[leftmargin=8mm,itemsep=1mm,topsep=2mm]
  \item Раскройте скобки.
  \item Перенесите слагаемые с \(x\) в одну сторону, числа --- в другую.
  \item Приведите подобные.
  \item Разделите на коэффициент при \(x\).
\end{enumerate}

При умножении или делении обеих частей неравенства на отрицательное число знак неравенства меняется на противоположный.

\exampletitle
\[
4x-1\ge0
\quad\Longrightarrow\quad
4x\ge1
\quad\Longrightarrow\quad
x\ge\frac14.
\]
Ответ:
\[
x\in\left[\frac14;+\infty\right).
\]

\exampletitle
\[
-4x-1\ge0
\quad\Longrightarrow\quad
-4x\ge1.
\]
Делим на \(-4\), поэтому меняем знак:
\[
x\le-\frac14.
\]

\important{Перед последним делением проверьте знак коэффициента при \(x\). Именно здесь чаще всего теряется смена знака неравенства.}

\exampletitle
\[
2(x-1)>-3(x-2)
\]
\[
2x-2>-3x+6
\quad\Longrightarrow\quad
5x>8
\quad\Longrightarrow\quad
x>\frac85.
\]

% =================================================
% 4. СИСТЕМЫ НЕРАВЕНСТВ
% =================================================
\section{4. Системы линейных неравенств}

\keyidea{Смысл системы}
Решение системы должно удовлетворять \textbf{каждому} неравенству одновременно. Поэтому после решения отдельных неравенств берём \textbf{пересечение} полученных множеств.

\exampletitle
Решим систему
\[
\begin{cases}
2x-1\ge0,\\
-3x+6\ge0.
\end{cases}
\]
Получаем
\[
x\ge\frac12,
\qquad
x\le2.
\]
Следовательно,
\[
x\in\left[\frac12;2\right].
\]

\begin{center}
\begin{tikzpicture}[x=1.55cm,y=1cm]
  \draw[-{Latex[length=2.5mm]},line width=0.8pt] (-0.5,0)--(5.5,0);

  \foreach \x/\lab in {1/{\frac12},4/{2}}{
    \draw (\x,0.08)--(\x,-0.08);
    \node[below=2mm] at (\x,0) {$\lab$};
  }

  \draw[accent,line width=1.8pt] (1,0.28)--(5.05,0.28);
  \filldraw[accent] (1,0.28) circle (1.6pt);

  \draw[darkaccent,line width=1.8pt] (-0.05,-0.28)--(4,-0.28);
  \filldraw[darkaccent] (4,-0.28) circle (1.6pt);

  \draw[textgray,line width=2.4pt] (1,0)--(4,0);
  \filldraw[textgray] (1,0) circle (1.6pt);
  \filldraw[textgray] (4,0) circle (1.6pt);
\end{tikzpicture}
\end{center}

\important{Если общей части нет, ответ системы --- \(\emptyset\). Например, условия \(x>3\) и \(x\le1\) одновременно выполнить невозможно.}

% =================================================
% FLOWCHART 2
% =================================================
\clearpage
\thispagestyle{plain}

\begin{center}
{\Large\bfseries\color{darkaccent}Алгоритм: система неравенств}
\end{center}

\vspace{7mm}

\begin{center}
\begin{tikzpicture}[node distance=11mm]

\node[flowstart] (start) {Дана система неравенств};

\node[flowaction, below=of start] (solve)
{Решить каждое неравенство отдельно};

\node[flowaction, below=of solve] (axis)
{Отметить решения на общей числовой прямой};

\node[flowdecision, below=14mm of axis] (cross)
{Есть общая часть?};

\node[flowside, below=18mm of cross, xshift=-52mm] (intersection)
{Выделить пересечение};

\node[flowside, below=18mm of cross, xshift=52mm] (empty)
{Записать \(\emptyset\)};

\node[flowstart, below=28mm of cross] (finish) {Ответ};

\draw[flowarrow] (start)--(solve);
\draw[flowarrow] (solve)--(axis);
\draw[flowarrow] (axis)--(cross);

\draw[flowarrow]
  (cross.west) -|
  node[flowlabel,pos=0.22]{да}
  (intersection.north);

\draw[flowarrow]
  (cross.east) -|
  node[flowlabel,pos=0.22]{нет}
  (empty.north);

\draw[flowarrow] (intersection.south) |- (finish.west);
\draw[flowarrow] (empty.south) |- (finish.east);

\end{tikzpicture}
\end{center}

\clearpage

% =================================================
% 5. СИСТЕМЫ УРАВНЕНИЙ
% =================================================
\section{5. Системы линейных уравнений: метод подстановки}

\keyidea{Идея метода}
Удобнее выражать ту переменную, у которой коэффициент равен \(1\) или \(-1\). Затем полученное выражение подставляют во второе уравнение.

\exampletitle
Решим
\[
\begin{cases}
2k+b=3,\\
-k+2b=4.
\end{cases}
\]
Из второго уравнения удобно выразить \(k\):
\[
-k+2b=4
\quad\Longrightarrow\quad
2b-4=k,
\]
то есть
\[
k=2b-4.
\]
Подставляем в первое уравнение:
\[
2(2b-4)+b=3,
\]
\[
4b-8+b=3,
\]
\[
5b=11,
\qquad
b=\frac{11}{5}.
\]
Тогда
\[
k=2\cdot\frac{11}{5}-4=\frac25.
\]
Ответ:
\[
(k;b)=\left(\frac25;\frac{11}{5}\right).
\]

\important{При переносе слагаемого через знак равенства его знак меняется. Если перед выражаемой переменной стоит \(-1\), иногда удобнее перенести всю эту переменную на другую сторону, чтобы сразу получить положительный коэффициент.}

% =================================================
% 6. РАЦИОНАЛЬНЫЕ УРАВНЕНИЯ
% =================================================
\section{6. Рациональные уравнения и общий знаменатель}

\keyidea{Перед преобразованиями}
Если переменная находится в знаменателе, сначала запишите ограничения: знаменатель не должен обращаться в ноль.

\exampletitle
Решим
\[
\frac{x+1}{x-1}-\frac{x-1}{x+1}=2.
\]
Ограничения:
\[
x\ne1,
\qquad
x\ne-1.
\]
Общий знаменатель:
\[
(x-1)(x+1)=x^2-1.
\]
После приведения к общему знаменателю:
\[
\frac{(x+1)^2-(x-1)^2}{x^2-1}=2.
\]
Используем разность квадратов:
\[
(x+1)^2-(x-1)^2
=
\bigl((x+1)-(x-1)\bigr)
\bigl((x+1)+(x-1)\bigr)
=4x.
\]
Поэтому
\[
\frac{4x}{x^2-1}=2.
\]
С учётом записанных ограничений умножаем обе части на \(x^2-1\):
\[
4x=2(x^2-1),
\]
\[
x^2-2x-1=0.
\]
Отсюда
\[
x=1\pm\sqrt2.
\]
Оба значения допустимы, так как не равны \(\pm1\).

Ответ:
\[
x=1-\sqrt2,
\qquad
x=1+\sqrt2.
\]

\important{При вычитании целого выражения заключайте его в скобки. Например, в выражении \((x+1)^2-(x-1)^2\) знак «минус» относится ко всему второму квадрату.}

\important{Нельзя делить на выражение, которое может быть равно нулю, или сокращать его без фиксации ограничения. После преобразований проверяйте найденные корни по ограничениям.}

% =================================================
% 7. РАСЧЁТЫ ПО ФОРМУЛЕ
% =================================================
\section{7. Расчёты по готовой формуле}

\keyidea{Что делать}
\begin{enumerate}[leftmargin=8mm,itemsep=1mm,topsep=2mm]
  \item Определите, какие величины даны в условии.
  \item Сопоставьте числа с буквами в формуле.
  \item Подставьте значения.
  \item Выполните вычисления и проверьте арифметику.
\end{enumerate}

\exampletitle
Если
\[
C=6000+400m
\]
и \(m=5\), то
\[
C=6000+400\cdot5=8000.
\]

\keyidea{Если нужно выразить величину}
Например, из
\[
S=\frac12 ab\sin\alpha
\]
при \(a\ne0\) и \(b\ne0\) получаем
\[
\sin\alpha=\frac{2S}{ab}.
\]
Промежуточные преобразования лучше записывать явно: так проще заметить потерянный множитель, знак или знаменатель.

% =================================================
% 8. МЕТОД ИНТЕРВАЛОВ
% =================================================
\section{8. Квадратные неравенства и метод интервалов}

\keyidea{Когда нужен метод интервалов}
Линейное неравенство можно привести к виду \(x>a\), \(x\le a\) и сразу получить ответ. В квадратном неравенстве, например
\[
x^2-x-6>0,
\]
переменная одновременно входит в \(x\) и \(x^2\), поэтому используем другой подход.

\subsection{Шаг 1. Приведите неравенство к стандартному виду}
Справа должен остаться ноль. Например,
\[
x^2-x>6
\quad\Longrightarrow\quad
x^2-x-6>0.
\]

\subsection{Шаг 2. Найдите нули левой части}
Временно решаем уравнение
\[
x^2-x-6=0.
\]
Его корни:
\[
x=-2,
\qquad
x=3.
\]
Они разбивают числовую прямую на интервалы
\[
(-\infty;-2),
\qquad
(-2;3),
\qquad
(3;+\infty).
\]

\subsection{Шаг 3. Определите знаки на интервалах}
Берём удобную контрольную точку из одного интервала. Для центрального интервала удобно взять \(x=0\):
\[
0^2-0-6=-6<0.
\]
Значит, на интервале \((-2;3)\) выражение отрицательно. Оба корня простые, поэтому при переходе через каждый из них знак меняется:
\[
+\qquad -\qquad +.
\]

\begin{center}
\begin{tikzpicture}[x=1.55cm,y=1cm]
  \draw[-{Latex[length=2.5mm]},line width=0.8pt] (-3.2,0)--(4.5,0);

  \filldraw[white,draw=textgray,line width=0.8pt] (-1.4,0) circle (2pt);
  \filldraw[white,draw=textgray,line width=0.8pt] (2.0,0) circle (2pt);

  \node[below=2mm] at (-1.4,0) {$-2$};
  \node[below=2mm] at (2.0,0) {$3$};

  \node[above=3mm] at (-2.35,0) {$+$};
  \node[above=3mm] at (0.3,0) {$-$};
  \node[above=3mm] at (3.05,0) {$+$};
\end{tikzpicture}
\end{center}

Поскольку требуется положительное значение левой части, выбираем интервалы со знаком «плюс»:
\[
x\in(-\infty;-2)\cup(3;+\infty).
\]

\important{Для строгого неравенства \(>0\) или \(<0\) нули выражения в ответ не входят. Для нестрогого \(\ge0\) или \(\le0\) допустимые нули включаются.}

\keyidea{Чётная и нечётная кратность корня}
Если множитель имеет вид
\[
(x-a)^m,
\]
то при переходе через \(x=a\):
\begin{itemize}[leftmargin=8mm,itemsep=1mm,topsep=2mm]
  \item при нечётном \(m\) знак меняется;
  \item при чётном \(m\) знак сохраняется.
\end{itemize}

Например, в выражении
\[
(x-2)^2(x+1)
\]
корень \(x=2\) имеет чётную кратность, а \(x=-1\) --- нечётную.

\important{После переноса удобно записывать многочлен по убыванию степеней: сначала \(x^2\), затем \(x\), затем свободный член. Это облегчает проверку коэффициентов и поиск корней.}

% =================================================
% FLOWCHART 3
% =================================================
\clearpage
\begingroup
\pdfpagewidth=297mm
\pdfpageheight=420mm
\paperwidth=297mm
\paperheight=420mm
\setlength{\textwidth}{247mm}
\setlength{\linewidth}{247mm}
\setlength{\columnwidth}{247mm}
\hsize=247mm
\setlength{\textheight}{376mm}
\setlength{\oddsidemargin}{-0.4mm}
\setlength{\evensidemargin}{-0.4mm}
\setlength{\topmargin}{-3.4mm}
\setlength{\headheight}{0pt}
\setlength{\headsep}{0pt}
\thispagestyle{empty}
\enlargethispage{120mm}

\begin{center}
{\Large\bfseries\color{darkaccent}Алгоритм: метод интервалов}
\end{center}

\vspace{6mm}

\begin{center}
\begin{tikzpicture}[node distance=9mm]

\node[flowstart] (start) {Дано неравенство};

\node[flowaction, below=of start] (normalize)
{Перенести все слагаемые влево, справа оставить \(0\)};

\node[flowaction, below=of normalize] (zeros)
{Решить уравнение «левая часть \(=0\)» и найти корни};

\node[flowaction, below=of zeros] (axis)
{Отметить корни на числовой прямой};

\node[flowaction, below=of axis] (sign)
{Определить знак на одном интервале контрольной подстановкой};

\node[flowdecision, below=13mm of sign] (mult)
{Кратность следующего корня нечётная?};

\node[flowside, below=17mm of mult, xshift=-65mm] (change)
{При переходе поменять знак};

\node[flowside, below=17mm of mult, xshift=65mm] (same)
{При переходе сохранить знак};

\node[flowaction, below=27mm of mult] (repeat)
{Повторить для остальных корней и отметить знаки всех интервалов};

\node[flowaction, below=of repeat] (select)
{Выбрать интервалы, соответствующие знаку исходного неравенства};

\node[flowdecision, below=13mm of select] (strict)
{Неравенство нестрогое?};

\node[flowside, below=17mm of strict, xshift=-65mm] (include)
{Включить допустимые граничные точки};

\node[flowside, below=17mm of strict, xshift=65mm] (exclude)
{Не включать граничные точки};

\node[flowstart, below=27mm of strict] (finish)
{Записать ответ};

\draw[flowarrow] (start)--(normalize);
\draw[flowarrow] (normalize)--(zeros);
\draw[flowarrow] (zeros)--(axis);
\draw[flowarrow] (axis)--(sign);
\draw[flowarrow] (sign)--(mult);

\draw[flowarrow]
  (mult.west) -|
  node[flowlabel,pos=0.2]{да}
  (change.north);

\draw[flowarrow]
  (mult.east) -|
  node[flowlabel,pos=0.2]{нет}
  (same.north);

\draw[flowarrow] (change.south) |- (repeat.west);
\draw[flowarrow] (same.south) |- (repeat.east);
\draw[flowarrow] (repeat)--(select);
\draw[flowarrow] (select)--(strict);

\draw[flowarrow]
  (strict.west) -|
  node[flowlabel,pos=0.2]{да}
  (include.north);

\draw[flowarrow]
  (strict.east) -|
  node[flowlabel,pos=0.2]{нет}
  (exclude.north);

\draw[flowarrow] (include.south) |- (finish.west);
\draw[flowarrow] (exclude.south) |- (finish.east);

\end{tikzpicture}
\end{center}

\clearpage
\endgroup

% =================================================
% 9. САМОПРОВЕРКА
% =================================================
\section{9. Короткая самопроверка}

\keyidea{Уравнения и преобразования}
\begin{itemize}[leftmargin=8mm,itemsep=1.5mm,topsep=2mm]
  \item Проверьте коэффициенты \(a\), \(b\), \(c\) и их знаки.
  \item При \(x^2=a\), \(a>0\), не теряйте знак \(\pm\).
  \item Ищите общий множитель и формулы сокращённого умножения.
  \item При переносе слагаемого меняйте его знак.
\end{itemize}

\keyidea{Дроби и ограничения}
\begin{itemize}[leftmargin=8mm,itemsep=1.5mm,topsep=2mm]
  \item До сокращения выпишите значения, обращающие знаменатель в ноль.
  \item Сокращайте только множители.
  \item Не теряйте ограничения после сокращения.
  \item При вычитании выражения ставьте скобки.
\end{itemize}

\keyidea{Неравенства}
\begin{itemize}[leftmargin=8mm,itemsep=1.5mm,topsep=2mm]
  \item При делении или умножении на отрицательное число меняйте знак неравенства.
  \item В системе берите пересечение решений.
  \item В методе интервалов выбирайте интервалы именно нужного знака.
  \item Учитывайте строгость неравенства и кратность корней.
\end{itemize}

\keyidea{Оформление решения}
Записывайте существенные промежуточные шаги. Они позволяют проверить, откуда появился каждый знак, множитель и знаменатель, и быстрее найти арифметическую ошибку.

% =================================================
% ТРЕНИРОВОЧНЫЙ БЛОК
% =================================================
\clearpage
\section{Тренировочный блок. Путь А}

\begin{center}
{\large\bfseries 7 заданий по нарастающей сложности}
\end{center}

\vspace{2mm}

\begin{enumerate}[leftmargin=9mm,label=\textbf{\arabic*.},itemsep=5mm]

\item Решите уравнение и разложите левую часть на множители:
\[
x^2-x-12=0.
\]

\item Стоимость услуги определяется формулой
\[
C=6000+400m.
\]
Найдите \(C\), если \(m=5\).

\item Упростите алгебраическую дробь и укажите ограничения:
\[
\frac{x^2-5x+6}{x^2-9}.
\]

\item Решите систему неравенств:
\[
\begin{cases}
3(2x+5)\ge -x+2,\\
4(x-3)<2(-x+2).
\end{cases}
\]

\item Решите систему методом подстановки:
\[
\begin{cases}
2x+y=7,\\
-x+2y=4.
\end{cases}
\]

\item Решите уравнение:
\[
\frac{x+2}{x-2}-\frac{x-2}{x+2}=\frac83.
\]

\item Решите методом интервалов:
\[
x^2-7x+10\le0.
\]

\end{enumerate}

\vfill

\noindent\textbf{Перед сдачей:} в заданиях с дробями проверьте ограничения; в неравенствах --- знак, граничные точки и выбранные интервалы.

% =================================================
% ОТВЕТЫ
% =================================================
\clearpage
\section{Ответы}

\renewcommand{\arraystretch}{1.45}

\begin{table}[ht]
\centering
\begin{tabularx}{\linewidth}{@{}cX@{}}
\toprule
\textbf{№} & \textbf{Ответ} \\
\midrule
1 & \(x=-3,\ 4;\quad x^2-x-12=(x+3)(x-4)\). \\
2 & \(C=8000\). \\
3 & \(\displaystyle \frac{x-2}{x+3},\quad x\ne-3,\ x\ne3\). \\
4 & \(\displaystyle x\in\left[-\frac{13}{7};\frac83\right)\). \\
5 & \(x=2,\ y=3\). \\
6 & \(x=-1,\ 4\). \\
7 & \(x\in[2;5]\). \\
\bottomrule
\end{tabularx}
\end{table}

\vfill

\begin{center}
{\small\color{textgray}\teacher\ эксклюзивно для \student}
\end{center}

\clearpage

\end{document}